% sec06_5.tex — Section 6.5: Wave Equation
\section{Wave Equation}
\label{sec:fd-wave}

We may also compute solutions to the one-dimensional wave equation by introducing finite difference approximations.  Using \textbf{centered differences in space and time}, the wave equation,
\begin{equation}
\boxed{\pdd{u}{t} = c^2\,\pdd{u}{x},}
\label{eq:6.5.1}
\end{equation}
becomes the following partial difference equation,
\begin{equation}
\boxed{\frac{u_j^{(m+1)} - 2u_j^{(m)} + u_j^{(m-1)}}{(\Delta t)^2} = c^2\,\frac{u_{j+1}^{(m)} - 2u_j^{(m)} + u_{j-1}^{(m)}}{(\Delta x)^2}.}
\label{eq:6.5.2}
\end{equation}
The truncation error is the sum of a term of $O(\Delta x)^2$ and one of $O(\Delta t)^2$.  By solving \eqref{eq:6.5.2} for $u_j^{(m+1)}$, the solution may be marched forward in time.  Three levels of time are involved in \eqref{eq:6.5.2}, as indicated in Fig.~\ref{fig:6.5.1}.  $u(x,t)$ is needed ``initially'' at two values of $t$ ($0$ and $-\Delta t$) to start the calculation.  We use the two initial conditions for the wave equation, $u(x,0) = f(x)$ and $\partial u/\partial t(x,0) = g(x)$, to compute at $m = 0$ and at $m = -1$.  Using a centered difference in time for $\partial u/\partial t$ [so as to maintain an $O(\Delta t)^2$ truncation error] yields
\begin{align}
u_j^{(0)} &= f(x_j) = f(j\Delta x) \label{eq:6.5.3} \\
\frac{u_j^{(1)} - u_j^{(-1)}}{2\Delta t} &= g(x_j) = g(j\Delta x). \label{eq:6.5.4}
\end{align}

\begin{figure}[H]
\centering
\includegraphics[width=0.8\textwidth]{figures/ch06/fig_6_5_1.png}
\caption{Marching forward in time for the wave equation.}
\label{fig:6.5.1}
\end{figure}

To begin the calculation we must compute $u_j^{(-1)}$.  The initial conditions, \eqref{eq:6.5.3} and \eqref{eq:6.5.4}, are two equations in three unknowns, $u_j^{(-1)}$, $u_j^{(0)}$, $u_j^{(1)}$.  The partial difference equation at $t = 0$ provides a third equation:
\begin{equation}
u_j^{(1)} = 2u_j^{(0)} - u_j^{(-1)} + \frac{c^2}{(\Delta x/\Delta t)^2}\left(u_{j-1}^{(0)} - 2u_j^{(0)} + u_{j+1}^{(0)}\right).
\label{eq:6.5.5}
\end{equation}
$u_j^{(1)}$ may be eliminated from \eqref{eq:6.5.4} and \eqref{eq:6.5.5} since $u_j^{(0)}$ is known from \eqref{eq:6.5.3}.  In this way we can solve for $u_j^{(-1)}$.  Once $u_j^{(-1)}$ and $u_j^{(0)}$ are known, the later values of $u$ may be computed via \eqref{eq:6.5.2}.  Boundary conditions may be analyzed as before.

\paragraph{Stability.}
Our limited experience should already suggest that a stability analysis is important.  To determine if any spatially periodic waves grow, we substitute
\begin{equation}
u_j^{(m)} = Q^{t/\Delta t}e^{i\alpha x}
\label{eq:6.5.6}
\end{equation}
into \eqref{eq:6.5.2}, yielding
\begin{equation}
Q - 2 + \frac{1}{Q} = \sigma,
\label{eq:6.5.7}
\end{equation}
where
\begin{equation}
\sigma = \frac{c^2}{(\Delta x/\Delta t)^2}\left(e^{i\alpha\Delta x} - 2 + e^{-i\alpha\Delta x}\right) = \frac{2c^2}{(\Delta x/\Delta t)^2}\left[\cos(\alpha\Delta x) - 1\right].
\label{eq:6.5.8}
\end{equation}
Equation~\eqref{eq:6.5.7} is a quadratic equation for $Q$ since \eqref{eq:6.5.2} involves three time levels and two time differences:
\begin{equation}
Q^2 - (\sigma + 2)\,Q + 1 = 0 \quad\text{and thus}\quad Q = \frac{\sigma + 2 \pm \sqrt{(\sigma+2)^2 - 4}}{2}.
\label{eq:6.5.9}
\end{equation}
The two roots correspond to two ways in which waves evolve in time.  If $-2 < \sigma + 2 < 2$, the roots are complex conjugates of each other.  In this case (as discussed earlier)
\[
Q^m = (re^{i\theta})^m = r^m e^{im\theta},
\]
where $r = |Q|$ and $\theta = \arg Q$.  Since
\[
|Q|^2 = \frac{(\sigma + 2)^2}{4} + \frac{4 - (\sigma+2)^2}{4} = 1,
\]
the solution oscillates for fixed $x$ as $m$ (time) increases if $-2 < \sigma + 2 < 2$.  This is similar to the wave equation itself, which permits time-periodic solutions when the spatial part is periodic (e.g., $\sin n\pi x/L\cos n\pi ct/L$).  If $\sigma + 2 > 2$ or $\sigma + 2 < -2$, the roots are real with product equal to~1 [see \eqref{eq:6.5.9}].  Thus, one root will be greater than one in absolute value, giving rise to an unstable behavior.

The solution will be stable if $-2 < \sigma + 2 < 2$ or $-4 < \sigma < 0$.  From \eqref{eq:6.5.8} we conclude that our numerical scheme is stable if
\begin{equation}
\boxed{\frac{c}{\Delta x/\Delta t} \le 1,}
\label{eq:6.5.10}
\end{equation}
known as the \textbf{Courant stability condition} (for the wave equation).  Here $c$ is the speed of propagation of signals for the wave equation and $\Delta x/\Delta t$ is the speed of propagation of signals for the discretization of the wave equation.  Thus, we conclude that for stability, \textbf{the numerical scheme must have a greater propagation speed than the wave equation itself}.  In this way the numerical scheme will be able to account for the real signals being propagated.  The stability condition again limits the size of the time step, in this case
\begin{equation}
\Delta t \le \frac{\Delta x}{c}.
\label{eq:6.5.11}
\end{equation}

\paragraph{Convergence.}
The time-dependent part oscillates if $-2 < \sigma + 2 < 2$ (the criterion for stability):
\begin{equation}
Q^{t/\Delta t} = e^{i\frac{t}{\Delta t}\arg Q}.
\label{eq:6.5.12}
\end{equation}
We note that if $\alpha\Delta x$ is small (where $\alpha$ is the wave number, usually $n\pi/L$), then $\cos(\alpha\Delta X) \approx 1 - \frac{1}{2}(\alpha\Delta x)^2$ from its Taylor series.  Thus, the parameter $\sigma$ is small and negative:
\[
\sigma \approx -\frac{1}{2}(\alpha\Delta x)^2\,\frac{2c^2}{(\Delta x/\Delta t)^2}.
\]
In this case the temporal frequency of the solutions of the partial difference equations is
\begin{equation}
\frac{\arg Q}{\Delta t} = \frac{\tan^{-1}\!\left(\frac{\sqrt{4 - (\sigma+2)^2}}{\sigma + 2}\right)}{\Delta t} \approx \frac{\left(\frac{\sqrt{4 - (\sigma+2)^2}}{\sigma + 2}\right)}{\Delta t} \approx \frac{(\sqrt{-\sigma})}{\Delta t} \approx c\alpha,
\label{eq:6.5.13}
\end{equation}
since $\sigma$ is very small and negative and $\tan\phi \approx \phi$ for small angles $\phi$.  Equation~\eqref{eq:6.5.13} shows that the temporal frequency of the partial difference equation approaches the frequency of the partial differential equation.  Perhaps this is clearer by remembering that $\alpha$ is the wave number usually $n\pi/L$.  Thus, if $\alpha\Delta x$ is small, the partial difference equation has solutions that are approximately the same as the partial differential equation:
\[
Q^{t/\Delta t}e^{i\alpha x} = e^{i\frac{t}{\Delta t}\arg Q}e^{i\alpha x} \approx e^{i\alpha x}e^{ic\alpha t}.
\]

It can also be shown that the error (difference between the solutions of the partial difference and differential equations at fixed $x$ and $t$) is $O(\Delta t)$ if $\alpha\Delta x$ is small.  It is usual to compute with fixed $\Delta x/\Delta t$.  To improve a computation, we cut $\Delta x$ in half (and thus $\Delta t$ will be decreased by a factor of~2).  Thus, if $\alpha\Delta x \ll 1$ with $\Delta x/\Delta t$ fixed, numerical errors (at fixed $x$ and $t$) are cut in half if the discretization step size $\Delta x$ is halved (and time step $\Delta t$ is halved).  All computations should be done satisfying the Courant stability condition \eqref{eq:6.5.10}.

\begin{exercises}{6.5}
\item Modify the Courant stability condition for the wave equation to account for the boundary condition $u(0) = 0$ and $u(L) = 0$.

\item Consider the wave equation subject to the initial conditions
\[
u(x,0) = \begin{cases} 1 & \frac{L}{4} < x < \frac{3L}{4} \\ 0 & \text{otherwise}\end{cases}
\]
\[
\pd{u}{t}(x,0) = 0
\]
and the boundary conditions
\begin{align*}
u(0,t) &= 0 \\
u(L,t) &= 0.
\end{align*}
Use nine interior mesh points and compute using centered differences in space and time.  Compare to the exact solution.
\begin{enumerate}
\item[(a)] $\Delta t = \Delta x/2c$ \qquad (b) $\Delta t = \Delta x/c$ \qquad (c) $\Delta t = 2\Delta x/c$
\end{enumerate}

\item For the wave equation, $u(x,t) = f(x - ct)$ is a solution, where $f$ is an arbitrary function.  If $c = \Delta x/\Delta t$, show that $u_j^m = f(x_j - ct_m)$ is a solution of \eqref{eq:6.5.2} for arbitrary $f$.

\item Show that the conclusion of Exercise~6.5.3 is not valid if $c \ne \Delta x/\Delta t$.

\item Consider the first-order wave equation
\[
\pd{u}{t} + c\,\pd{u}{x} = 0.
\]
\begin{enumerate}
\item[(a)] Determine a partial difference equation by using a forward difference in time and a centered difference in space.
\item[\starred (b)] Analyze the stability of this scheme (without boundary conditions).
\end{enumerate}

\item[\starred 6.5.6.] Modify Exercise~6.5.5 for centered difference in space and time.

\item Solve on a computer [using \eqref{eq:6.5.2}] the wave equation $\partial^2 u/\partial t^2 = \partial^2 u/\partial x^2$ with $u(0,t) = 0$, $u(1,t) = 0$, $u(x,0) = \sin\pi x$ and $\partial u/\partial t(x,0) = 0$ with $\Delta x = 1/100$.  For the initial condition for the first derivative $\partial u/\partial t(x,0) = 0$, use a forward difference instead of the centered difference \eqref{eq:6.5.4}.  Compare to analytic solution at $x = 1/2$, $t = 1$ by computing the error there (the difference between the analytic solution and the numerical solution).  Pick $\Delta t$ so that
\begin{enumerate}
\item[(a)] $\Delta x/\Delta t = 1.5$ \qquad (b) $\Delta x/\Delta t = 0.5$
\item[(c)] To improve the calculation in part~(a), let $\Delta x = 1/200$ but keep $\Delta x/\Delta t = 1.5$.
\item[(d)] Compare the errors in parts~(a) and~(c).
\end{enumerate}
\end{exercises}
